\documentclass[11pt]{article}
\usepackage[utf8]{inputenc}
\usepackage[margin=1in]{geometry}
\usepackage{amsmath,amssymb,amsfonts}
\usepackage{booktabs}
\usepackage{microtype}
\usepackage{hyperref}
\usepackage{titlesec}
\usepackage{enumitem}

\titleformat{\section}{\large\bfseries}{\thesection}{1em}{}[{\titlerule[0.5pt]}]
\titleformat{\subsection}{\normalsize\bfseries}{\thesubsection}{1em}{}

\title{\textbf{Executive Research Summary: Graph Digital Twin for Smart Grid Cyber-Resilience}}
\author{\textbf{Research Project Overview \& Novelty Analysis} \\ Submitted for Faculty Advisor Review}
\date{\today}

\begin{document}

\maketitle

\section{Project Objectives}
This project develops an explainable, topology-invariant \textbf{Graph Digital Twin (GDT)} framework to protect smart grids from stealthy False Data Injection Attacks (FDIAs) targeting SCADA telemetry. By treating the physical electricity network as a dynamic spatial graph, our system concurrently detects grid-level tampering and localizes compromised substation sensors in a single, real-time forward pass.

\section{Base Literature \& Context}
Our framework builds directly upon and extends two recent publications:
\begin{enumerate}[label={[\arabic*]}]
    \item \textbf{Base Paper:} S. Araar, B. Farou, Z. Kouahla, and H. Seridi, ``Topology-Aware Cyberattack Detection in Smart Grids Using Explainable Graph Neural Networks,'' in \textit{Proc. Third National Conference on Applications of Artificial Intelligence (A2I-26)}, Boumerdes, Algeria, 2026.
    \item \textbf{2nd Reference Paper:} X. Yue and G. Qu, ``Graph neural networks approach for cyberattack detection in smart grids,'' in \textit{Proc. IEEE International Conference on Electro Information Technology (EIT)}, 2024, pp. 1--6.
\end{enumerate}

\section{Research Gaps and Proposed Novelties}
Prior GNN-based grid security works exhibit critical gaps that limit real-world deployment. The table below outlines how our proposed Graph Digital Twin resolves these limitations:

\begin{table}[h!]
\centering
\small
\renewcommand{\arraystretch}{1.3}
\begin{tabular}{p{3.2cm}p{6.2cm}p{6.2cm}}
\toprule
\textbf{Focus Area} & \textbf{Prior Art Gaps (Araar [1], Yue [2])} & \textbf{Proposed GDT Solutions (Our Work)} \\
\midrule
\textbf{Task Complexity} & Limited to grid-level binary detection (passive alarm only). They do not locate the compromised target. & \textbf{Joint Multi-Task Architecture:} Concurrently detects attacks (99.36\% F1) and localizes substation anomalies (99.03\% F1). \\
\textbf{Grid Healing} & Passive indicators; state estimation remains corrupted, risking wide-area outages. & \textbf{Closed-Loop Mitigation:} Isolates flagged telemetry and substitutes baseline forecasts to recover \textbf{35.20\%} of Active Power MSE. \\
\textbf{Explainability Latency} & Uses GNNExplainer [1] (requires mask optimization per sample), which is too slow for real-time SCADA. & \textbf{Axiomatic Feature Attribution:} Deploys \textbf{Integrated Gradients} (gradient-based backpropagation) for instant attributions. \\
\textbf{Topology Shift} & Assumes static layout; model accuracy breaks when lines trip. & \textbf{Topology-Invariance:} Proved zero-shot generalization under random outages (maintains 70.59\% F1). \\
\bottomrule
\end{tabular}
\end{table}

\section{Core Experimental Results (IEEE 14-Bus Feeder)}
\begin{itemize}[leftmargin=*]
    \item \textbf{Nominal Performance:} The GraphSAGE backbone achieves a \textbf{99.36\% Detection F1-score} and a \textbf{99.03\% Localization F1-score}, outperforming standard GCNs by 25.19\% on localization due to its concat design.
    \item \textbf{Noise Sensitivity Boundary:} Under high Gaussian measurement noise ($\sigma = 0.08$), the pre-trained model maintains a robust \textbf{95.15\% Detection F1} and \textbf{87.76\% Localization F1}.
    \item \textbf{Grid Reconfiguration Resilience:} Under random single-line outages (dynamic grid topology reconfigurations), our pre-trained GNN generalizes robustly without retraining, whereas traditional flat MLPs fail (localization F1 drops to \textbf{45.78\%}) due to input dimensions mismatch.
\end{itemize}

\end{document}
